\section{Bounds: partial identification of \texorpdfstring{$q$}{q}}\label{sec:bounds}

This section delivers a sharp identified set for $q$ that any source-aggregation rule must respect. We begin with a benchmark assumption, then weaken it. Proofs are in Appendix~\ref{app:proofs}.

\begin{assumption}[Exchangeable collateral]\label{ass:ec}
Civilians of all demographic classes face equal per-capita death hazard.
\end{assumption}

Under Assumption~\ref{ass:ec} the combatant share is point-identified by a single rearrangement: civilian deaths split between the women-and-children class and the adult-male class in proportion to their civilian population shares, so
\begin{equation}\label{eq:ident}
q = 1 - \frac{\omega(1 - f)}{w}.
\end{equation}
The intuition is direct: $\omega$ falls short of the population share $w$ only to the extent that deaths are diverted to the (all-male) combatant pool.

\begin{remark}\label{rem:female-combatants}
Suppose a fraction $\eta\le\bar\eta=0.03$ of combatant \emph{deaths} is recorded in class $W$. This is a classification slack: the combatant stock itself stays in $\mathrm{AM}$, so the civilian split $S=w/(w+\mu(a-f))$ is unchanged. The identity becomes $\omega=(1-q)S+\eta q$, the true share is $q_{\mathrm{true}}=(S-\omega)/(S-\eta)$, and the no-slack value $q_{\eta=0}=(S-\omega)/S$ understates it by exactly
\[
q_{\mathrm{true}}-q_{\eta=0}\;=\;\frac{\eta\,q_{\mathrm{true}}}{S}\;=\;\frac{\eta\,q_{\eta=0}}{S-\eta}\;\ge\;0
\qquad(S>\eta,\ q_{\mathrm{true}}\ge 0).
\]
At the values in Section~\ref{sec:gaza}, where $S>\bar\eta$ throughout ($S\ge 0.597$ for $\mu\le 2$) and $q_{\eta=0}\le 0.26$, a loose uniform bound over $\mu\in[1,2]$ is $\bar\eta\cdot 0.26/(0.597-\bar\eta)\approx 1.4$ percentage points; the exact understatement is smaller and falls with $\mu$ ($1.05$ points at $\mu=1$, $0.33$ at $\mu=2$), since $q_{\eta=0}$ shrinks faster than $S$. Readers who want this slack should widen the \emph{upper} endpoints below by $\bar\eta\, q_{\eta=0}/(S-\bar\eta)$; we otherwise suppress it.
\end{remark}

Assumption~\ref{ass:ec} is strong: civilian adult men typically face higher per-capita exposure than women and children due to outside work, food queues, rescue activity, and circumstantial proximity to combat. We weaken it to a one-parameter family.

\begin{assumption}[$\mu$-bounded differential exposure]\label{ass:mu}
There exists $\mu\in[\mu_{\mathrm{lo}},\mu_{\mathrm{hi}}]$ such that, conditional on being civilian, class $\mathrm{AM}$ has per-capita death rate $\mu$ times that of class $W$.
\end{assumption}

\begin{theorem}[Sharp identified set]\label{thm:bounds}
Suppose $f\le a$ and Assumption~\ref{ass:mu} holds. Given $(\omega,w,a,f)$ and the exposure range $[\mu_{\mathrm{lo}},\mu_{\mathrm{hi}}]$, the identified set for $q$ is the image of the monotone-decreasing curve
\begin{equation}\label{eq:qmu}
q(\mu) = 1 - \omega\,\frac{w+\mu(a-f)}{w}, \qquad \mu\in[\mu_{\mathrm{lo}},\mu_{\mathrm{hi}}],
\end{equation}
giving the sharp interval $[q(\mu_{\mathrm{hi}}),\,q(\mu_{\mathrm{lo}})]\cap[0,1]$.
\end{theorem}

\begin{corollary}[Manpower bound]\label{cor:popbudget}
$q\le M/D$, so whenever the endpoints below are ordered, the joint feasible interval is
\[
q \in \bigl[\max\{0,\,q(\mu_{\mathrm{hi}})\},\; \min\{1,\,q(\mu_{\mathrm{lo}}),\,M/D\}\bigr];
\]
if the lower endpoint exceeds the upper the set is empty and the inputs are jointly infeasible---the event detected by a positive contradiction radius in Section~\ref{sec:contradiction}.
\end{corollary}

\paragraph{Calibrating the exposure range.} The exposure ratio $\mu$ is the framework's only behavioural parameter, and it can be disciplined empirically rather than stipulated. \citet{frost2026} identifies the civilian male exposure differential from the male-to-female death ratio just \emph{above} the combatant age range---where the combatant share is mechanically near zero---and validates the procedure on historical Gaza conflicts recorded by B'tselem, where combatant status is known: there, the observed male-to-female ratio at the top of the combatant age band is $3.4{:}1$ while the effective differential through the combatant ages is $5{:}1$, so the observable ratio is itself a conservative lower bound. Airstrike-specific data suggest smaller differentials ($\approx 1.3$; \citealp{cockerill2024}), consistent with indiscriminate mechanisms; ground operations push the differential up. We therefore report the identified set on a grid $\mu\in\{1, 1.5, 2, 2.5, 3.5, 5\}$ throughout, treating $[2,3.5]$ as the empirically best-supported range for a mixed air-and-ground campaign and the extreme values as stress tests. Larger $\mu$ \emph{lowers} the implied combatant share; the reader who disputes a high-$\mu$ calibration is pushed toward \emph{higher} $q$, and the sensitivity grid (Table~\ref{tab:sensitivity}) makes the entire mapping explicit.

\paragraph{Plug-in estimator and standard error.} The sample analogue of the upper endpoint is $\hat q=1-\hat\omega(1-\hat f)/\hat w$, with delta-method variance
\[
\widehat{\mathrm{Var}}(\hat q) \approx
\frac{\hat\omega^2}{\hat w^2}\widehat{\mathrm{Var}}(\hat f)
+\frac{(1-\hat f)^2}{\hat w^2}\widehat{\mathrm{Var}}(\hat\omega)
+\frac{\hat\omega^2(1-\hat f)^2}{\hat w^4}\widehat{\mathrm{Var}}(\hat w).
\]
